\subsection{Return Semantics and Frame Termination}

A function call creates a temporary execution frame containing local bindings and the current instruction pointer. When the function finishes, the frame terminates and the caller receives one object reference---explicitly via \texttt{return}, implicitly as \texttt{None}, or as a tuple when \texttt{return a, b} packs multiple values.

\subsubsection{Execution Frames: Inside the CPython Call Stack and \texttt{PyFrameObject} Allocation}

Unlike C, where activation records typically live on the hardware stack, CPython allocates evaluation frames as discrete \texttt{PyFrameObject} structures on the OS heap. Each frame holds \texttt{f\_code}, \texttt{f\_locals}, \texttt{f\_globals}, \texttt{f\_lasti} (instruction offset), and a link \texttt{f\_back} to the caller frame.

This design trades stack efficiency for runtime flexibility: frames can be introspected via \texttt{inspect.currentframe()}, preserved across \texttt{yield} suspensions (\texttt{gi\_frame}), and chained during recursion. The cost is heap allocation and reference management on every call---significant compared to bare-metal stack frames, but necessary for Python's object model and generator/coroutine semantics.

%<<<PROGRAM:programs/ret04>>>


Recursive calls create a chain of simultaneous heap frames, each running the same \texttt{co\_code} with different local state, until \texttt{sys.getrecursionlimit()} is exceeded.

\subsubsection{Implicit Return Defaults: Why Functions Without \texttt{return} Produce \texttt{None}}

A function that reaches the end of its body without \texttt{return} yields \texttt{None}. \texttt{print()} sends text to stdout; \texttt{return} sends an object to the caller. These are separate operations.

%<<<PROGRAM:programs/ret03>>>


An explicit bare \texttt{return} also returns \texttt{None}. There is normally one \texttt{None} object per process; check with \texttt{is None}.

\subsubsection{Multiple Return Values as Tuple Packing: The Real Structure Behind \texttt{return a, b}}

Python functions always return one object reference. \texttt{return a, b} constructs a tuple and returns it; the caller may unpack with \texttt{x, y = func()}.

%<<<PROGRAM:programs/ret02>>>


Starred unpacking (\texttt{first, *rest = func()}) collects overflow into a new \texttt{list}. Mismatch between target count and tuple length raises \texttt{ValueError} in the caller, not in the callee.

\subsubsection{Recursive Calls: Repeated Frame Allocation, Base Cases, and Recursion Depth Boundaries}

Each recursive call allocates a new heap frame with its own locals. The base case stops frame creation; returns unwind one object at a time.

%<<<PROGRAM:programs/ret01>>>


Without a base case or progress toward it, Python raises \texttt{RecursionError: maximum recursion depth exceeded}. Each active call consumes stack depth because every frame must remain allocated until its callee returns.
